\begin{savequote}
\includegraphics[width=3cm]{./images/teaser/affe-sama-ambas-schein.jpg}
\end{savequote}
\chapter
[\CO{[DOK]} Dokumentation]
{\CC{[DOK]} Dokumentation}
\label{chp:EIN}

\ChapterIntro
{%
Die medizinische Dokumentation
muss eine Geschichte erzählen. Eine Geschichte die ein
Dritter, unbeteiligter, verstehen \E{muss}. Dieser
unbeteiligte Dritte soll allein anhand der
Schilderung in Ihrer Dokumentation fähig sein eine
Verdachtsdiagnose zu stellen. Weiters sollte er
nachvollziehen können, was mit dem Patient passiert ist.
Auch die Ereignisse davor oder die Einsatzsituation soll
für denjenigen verständlich sein (Unfallhergang,
Wohnsituation bei psychiatrischen Patienten, \ldots)
}
{%
\begin{description}
  \item[Maintainer:] --
  \item[Autoren:] Diverse
  \item[Reviewer:] Standard-Reviewprozess
  \item[Version:] {\DocVersion} ({\DocVersionInfo})
  \item[SHA1:] {\DocGitCommit}
\end{description}

{\TxTeilkapitelAASS}
}



\clearpage % TEMP temp clearpage

\section{Allgemeines zur Dokumentation}

\index{Dokumentation!medizinische Aspekte}
\label{topic:dokumentation-medizinisch}

\Layout{0}{\TxQuerverweise}
{%
\begin{itemize}
  \item Rechtliche Aspekte:
  \FullPrettyRef{topic:dokumentation-rechtlich},
  \FullPrettyRef{topic:dokumentationspflicht-fallbeispiel}
  \item Dokumentation anderer Gesundheitsberufe (Anamnese):
  \FullPrettyRef{topic:dokumentation-andere-gesundheitsberufe}
\end{itemize}
}{%
\begin{itemize}
  \item Rechtlich:
  \prettyref{topic:dokumentation-rechtlich},
  \prettyref{topic:dokumentationspflicht-fallbeispiel}
  \item Andere Gesundheitsberufe:
  \prettyref{topic:dokumentation-andere-gesundheitsberufe}
\end{itemize}
}{}{}{}

\Layout{0}{Dokumentation -- Oder: Damit sich auch ein
Anderer auskennt \ldots} {Die medizinische Dokumentation
muss eine Geschichte erzählen. Eine Geschichte die ein
Dritter, unbeteiligter, verstehen \E{muss}. Dieser
unbeteiligte Dritte soll allein anhand der
Schilderung in Ihrer Dokumentation fähig sein eine
Verdachtsdiagnose zu stellen. Weiters sollte er
nachvollziehen können, was mit dem Patient passiert ist.
Auch die Ereignisse davor oder die Einsatzsituation soll
für denjenigen verständlich sein (Unfallhergang,
Wohnsituation bei psychiatrischen Patienten, \ldots)
}{
}
{}
{}
{Gabriel}

\bigskip

\Fazit{Die Dokumentation dient sowohl der rechtlichen
Absicherung, als auch der Beschreibung der Situation
gegenüber Dritter, welche am Einsatz nicht
beteiligt waren.}


\Layout{2}
{Dokumentationspflicht}
{\label{topic:dokumentationspflicht}%
\label{topic:dokumentation-rechtlich}%
Schon nach dem Sanitätergesetz besteht für jeden Sanitäter
die Pflicht, die von ihm gesetzten
sanitätsdienstlichen Maßnahmen zu
dokumentieren.\ZFootnote{{\S}\,5 SanG} Darüber hinaus müssen
analog zu anderen Gesundheitsberufen auch alle anderen
medizinisch relevanten Sachverhalte schriftlich und
nachvollziehbar fixiert werden.

Hintergrund dieser Dokumentationspflicht ist der
\EE{Heilbehandlungsvertrag}.
Im österreichischen Recht hat grundsätzlich der Kläger den
Schaden, die Rechtswidrigkeit und das Verschulden des
Beklagten zu beweisen. Besteht aber zwischen Kläger und
Beklagten ein Vertrag, so wird gem. \Z{{\S}\,1298} ABGB das
Verschulden des Beklagten angenommen bzw. der Beklagte muss
beweisen, dass ihn kein Verschulden trifft
(\T{Beweislastumkehr}).
Um nun für den Beklagten die Möglichkeit zu schaffen, sein
Unverschulden zu beweisen, wurde die
\I{Dokumentationspflicht} eingeführt, mit der bei
ausreichender Dokumentation wieder der Kläger das
Verschulden des Beklagten zu beweisen hat.


Inhalt der Dokumentation sollte nicht nur die 
gesetzten Maßnahmen sein, sondern auch 
\begin{itemize}
  \item situationsrelevante
  Umstände (beim Patient; am BO; während des Transportes;
  {\dots}),
  \item  u.\,U. relevante Inhalte einer Aufklärung,
  \item Einwilligung bzw. Ablehnung durch den Patienten
  (Revers) sowie
  \item sonstige für den Sanitäter entscheidend wirkende
  Umstände.
\end{itemize}



Es muss \E{gründlich}, \E{ausführlich} und
\E{nachvollziehbar} dokumentiert werden.
Die Dokumentation soll vorrangig am Einsatzprotokoll selbst
oder, wenn nötig, auf einem beigelegten
Blatt Papier erfolgen.
}
{}
{}
{}
{}


% \clearpage % TEMP temp clearpage
\section{Beispiele}

\Layout{2}{Schlechte Dokumentation}
{\mbox{}\par\bigskip\noindent\begin{minipage}{\linewidth}
{\sffamily Diagnose:} {\ttfamily\color{darkBlue} RQW}

{\sffamily Befunde:} {\sffamily RR}
{\FontTypewriter\color{darkBlue} ~~130/70,} {\sffamily HF}
{\FontTypewriter\color{darkBlue} ~~~90} {\sffamily SpO2}
{\FontTypewriter\color{darkBlue} ~~~98\% }

{\sffamily Anmerkungen: }{\FontTypewriter\color{darkBlue} -- --
--}
\end{minipage}}
{}
{}
{}
{}


\Layout{2}{Bessere Dokumentation}
{
\mbox{}\par\bigskip\noindent\begin{minipage}{\linewidth}
{\sffamily Diagnose:} {\FontTypewriter\color{darkBlue} RQW Stirn}

{\sffamily Befunde:} {\sffamily RR}
{\FontTypewriter\color{darkBlue} ~~130/70,} {\sffamily HF}
{\FontTypewriter\color{darkBlue} ~~~90} {\sffamily SpO2}
{\FontTypewriter\color{darkBlue} ~~~98\% }

{\sffamily Anmerkungen:} {\FontTypewriter\color{darkBlue} Pat.
fuhr mit Fahrrad und prallte mit Stirn gegen Ast., kein
Sturz. Lt. Begleiter keine B.losigkeit, Unfallhergang
erinnerlich \underbar{Traumacheck:} RQW Stirn, ca 50\,ml
Blutverlust sonst unauffällig. \underbar{Neuro:} grob
neurologisch unauffällig, Bew. klar, Pat zeitl, örtl, zur
Situation \& Person orientiert, Pupillen mittelweit, Reakt.
prompt und seitengleich, Kraft\&Gefühl seitengl. in OE und
UE. \underbar{Anamn.:} \underbar{Sympt:} Kopfschmerz seit
Aufprall. \underbar{Krankh.:} keine bek. \underbar{All.}
keine bek. \underbar{Med.:} Aspirin heute früh. letzte
Tetanus-Impf. nicht bek. \underbar{Ereignisse} s.o.
\underbar{Letzte} Mahlzeit: heute Früh 2 Semmeln\A{ KOCH:
? muss sein? GABS: ja, 1) wegen s-kamel, 2) wegen nüchtern }

\underline{Maßnahmen:} Steriler Verband, nüchtern belassen,
Transport liegend mit erh OK.}
\end{minipage}
}
{}
{}
{}
{}



\Layout{2}{Psychiatrischer und psychosozialer Patient}
{\mbox{}\par\bigskip\noindent\begin{minipage}{\linewidth}
{\sffamily Diagnose:} {\FontTypewriter\color{darkBlue}
Psychiatrische Erkrankung, V.a. Wahnvorstellungen, Unterbringung nach UbG}

{\sffamily Befunde:} {\sffamily RR}
{\FontTypewriter\color{darkBlue} ~~130/70,} {\sffamily HF}
{\FontTypewriter\color{darkBlue} ~~~90} {\sffamily SpO2}
{\FontTypewriter\color{darkBlue} ~~~98\% }

{\sffamily Anmerkungen:} {\FontTypewriter\color{darkBlue} Pat.
sitzend in Wohnung angetr., von Pol. mit Handschellen gefesselt. Pat. warf lt. Polizei Gegstände aus d Fenster
auf Straße, verletzte Passanten ($\rightarrow$ FLO-1). Pat.
gibt an, sich geg. Geheimdienst gewehrt zu haben,
jetzt agitiert und unkooperativ. 

\underbar{Traumacheck:} keine äußeren Anz f Verletzungen,
sonst nicht erhebbar (unkoop.). \underbar{Neuro:} Pat.
ist wach, sonst nicht erhebbar

\underbar{Anamn.:} \underbar{Sympt:} keine Antwort
\underbar{Krankh.:} lt. Spitalbrief Otto-Wagner-Sp. v.
1.4.08 paranoide Schizophrenie \underbar{All.} k.A.
\underbar{Med.:} nicht eruierbar, auf Tisch mehrere
angebr. Packungen Rohypnol \underbar{Ereignisse} s.o.
\underbar{Letzter Spitalsaufenthalt} k.A.

\underbar{Wohnsituation:} heruntergekommen, vermüllt, Pat.
wohnt mit großem Hund, ganze Wohnung verkotet, bestialisch
stinkend. 

\underbar{Maßnahmen:} Einweisung nach UbG, Transport in
Handschellen in Pol-begl. (Viktor 6, DNr. 4711). Protokoll
Amtsarzt ad. Spital.}
\end{minipage}}
{}
{}
{}
{}





% \section{Reserviert}
% \A{Medizinische Fachbereiche}
% 
% \section{Reserviert}
% \A{Medizinische Berufsgruppen}
% 
% \section{Reserviert}

